\chapter{Discussion}
\label{ch:discussion}

\section{Empirical Implications}
\label{sec:discussion-empirical}

The empirical results in Chapter~\ref{ch:results} support three claims that
matter for the operational deployment of LLM-based explanation systems on
top of formally verified ICS detectors.

\paragraph{ReasonGuard discriminates between local models.}
On the same prompt set, the clean-response rate (no V1--V5 fired) varies
substantially across the three local models tested. The framework therefore
acts as a model-selection signal: an organisation can compare candidate
models on the same bound set and pick the one with the best clean rate
under their bound configuration, rather than relying on aggregate benchmark
scores that may not reflect the safety-critical setting.

\paragraph{Violation types are not uniform across event types.}
Per-event-type analysis shows that all three local models produce a higher
violation rate on rarer proxy event types (network\_scan and
function\_code\_violation) than on the dominant normal\_modbus\_polling
class. This is consistent with the hypothesis that LLMs lean on
training-distribution priors when the formal bound is sparse: with a small
number of events of a given type in the bound, the model has fewer cues to
anchor its explanation on the formal evidence, and reaches for plausible
but unsupported claims.

\paragraph{The spaCy refinement matters in practice.}
The W11 refinement replaced the multi-pattern substring negation list with
a spaCy dependency-parse filter. On the W12 batch, V1 firings dropped by
approximately two-thirds on \texttt{llama3.1:8b-instruct-q4\_0} and by
approximately one-half on \texttt{phi3:mini}. This is a methodological
finding in its own right: lexicon-based detectors materially benefit from a
syntactic-scope layer when applied to real LLM outputs, because real LLM
outputs phrase negation in ways that hand-curated patterns rarely cover.

\section{Threats to Validity}
\label{sec:discussion-threats}

\paragraph{Construct validity.}
\textsc{ReasonGuard} treats the formal verdict as a bound on what an
explanation may correctly claim. The bound is therefore as strong as the
formal verdict it carries. Under the proxy regime
(\texttt{schema\_status: proxy\_until\_official\_automaton\_schema}) the
formal\_class field is generated by an internal classifier and is labelled
as a proxy throughout every artefact; the empirical conclusions in
Chapter~\ref{ch:results} are therefore conditional on the proxy bound and
will be repeated against the official schema once it is integrated.

\paragraph{Internal validity.}
The V1--V5 detector is deterministic given a fixed bound and a fixed
explanation. The main internal-validity concern is the manual annotation
itself: a single annotator scores all 200 rows, which introduces annotator
bias. The inter-rater agreement with the VUT Brno domain expert
(Section~\ref{sec:results-kappa}) is the headline mitigation; the M5
milestone defines $\kappa > 0.70$ as the target.

\paragraph{External validity.}
The evaluation is conducted on IEC-104 traffic from VUT Brno testbeds. The
framework is protocol-agnostic by construction --- nothing in the bound
schema or the verification mechanism is IEC-104-specific --- but
generalisation to Modbus, DNP3, S7, or non-ICS settings is not
empirically claimed in this thesis.

\paragraph{Statistical validity.}
The current local-model batch size (100 prompts per model under stratified
sampling) is sufficient to detect large per-model differences but
underpowered for fine-grained per-event-type comparisons within a single
model. The thesis reports the observed counts together with the relevant
sample sizes; statistical significance tests are reported only where the
sample size is sufficient for them to be meaningful.

\section{Operational Deployment Guidelines}
\label{sec:discussion-deployment}

The empirical findings translate into the following deployment guidelines
for organisations that intend to use an LLM as an explanation layer on top
of a formally verified ICS detector. The guidelines are written as a
checklist; each item is mapped to the EU AI Act provision it addresses in
Section~\ref{sec:discussion-eu-ai-act}.

\begin{enumerate}
  \item \textbf{Gate high-severity explanations.} Every explanation tagged
        by \textsc{ReasonGuard} as V1, V2, or V5 should be withheld from
        the operator and routed to a manual review queue. The operational
        workflow should not show high-severity explanations as
        authoritative.
  \item \textbf{Disclose the formal bound.} The operator should be able to
        inspect the bound that generated the explanation. The Streamlit
        dashboard in Section~\ref{sec:implementation-notes} is the
        canonical form of this disclosure.
  \item \textbf{Log every (bound, explanation, verdict) tuple.} The MLflow
        store at \texttt{src/pipeline\_config.MLFLOW\_TRACKING\_URI} is the
        audit trail; rolling backups of the SQLite file should be kept for
        the retention period required by the regulatory regime.
  \item \textbf{Re-evaluate when the underlying model changes.} Every
        model version, quantisation level, or prompt revision invalidates
        the prior evaluation. The pipeline supports tagging an evaluation
        run with the model version and the bound-schema version; the
        re-evaluation should be performed before the new artefact reaches
        production.
  \item \textbf{Communicate uncertainty explicitly.} Under the proxy
        schema, the explanation must not state a confirmed attack or
        normality conclusion. The forbidden-claim list in the bound makes
        this enforceable; the deployment workflow should additionally
        surface the schema status to the operator.
\end{enumerate}

\section{Mapping to the European Union AI Act}
\label{sec:discussion-eu-ai-act}

\subsection{Article 13 — Transparency and Information to Deployers}
\label{subsec:eu-article-13}
Article 13 requires high-risk AI systems to be transparent enough for the
deployer to interpret system output. The V1--V5 verdict directly addresses
this: it reports not just whether the explanation is unfaithful but
\emph{which} kind of unfaithfulness occurred and \emph{which} claim
triggered it. Deployment Guidelines (1) and (2) above operationalise the
transparency requirement.

\subsection{Article 17 — Quality Management System}
\label{subsec:eu-article-17}
Article 17 requires high-risk AI providers to maintain a quality
management system covering, among other things, evaluation, monitoring,
and corrective action. The MLflow-tracked evaluation pipeline plus the
manual annotation workflow plus the W13--14 F1 / $\kappa$ targets
collectively constitute such a system for the explanation layer.
Deployment Guidelines (3) and (4) operationalise the quality-management
requirement.

\subsection{Annex III — Critical Infrastructure}
\label{subsec:eu-annex-iii}
Annex III lists critical-infrastructure operation among the high-risk
categories. The severity-based gating in Deployment Guideline (1) maps
directly to the risk-management requirement that high-impact outputs be
subject to additional review before they affect operational decisions.

\section{Limitations and Future Work}
\label{sec:discussion-future}

The principal limitations of the current implementation are:

\paragraph{Proxy schema.} The bound's \texttt{formal\_class} field is
generated by a proxy classifier. The empirical findings are therefore
conditional on the proxy. The integration of the official NES@FIT
automaton schema is the highest-priority future work.

\paragraph{Rule-based classifier.} The V1--V5 classifier is rule-based. A
learning-based classifier trained on the 200-pair annotated set is a
natural extension; the manual annotation already produces the training
data, and the model versioning in MLflow already provides the evaluation
scaffolding.

\paragraph{Limited model set.} Three local and two cloud models are not
sufficient to draw general conclusions about LLM-class behaviour on this
task. Expanding to more models, including instruction-tuned variants of
the same base model, is straightforward in the current pipeline.

\paragraph{No operator-in-the-loop study.} The thesis does not measure
whether the V1--V5 severity classification changes operator behaviour
relative to the unguarded baseline. An operator study is essential before
the framework can be considered ready for production deployment, and is
proposed as the primary follow-up.
